\part{Le déploiement}
\label{part:deploiement}

\chapter{Monter le parc avec Vagrant}

\section{Quatre machines décrites en un fichier}

\begin{lstlisting}[language=Ruby,caption={\fichier{Vagrantfile} --- le parc, paramétrable}]
# Trois workers, comme le demande le cahier des charges. Sur un poste modeste,
# WORKERS=2 (ou 1) reduit le parc sans toucher au fichier :
#   WORKERS=2 vagrant up
NB_WORKERS = (ENV["WORKERS"] || "3").to_i.clamp(1, 9)

MACHINES = {
  "controller" => { "ip" => "192.168.56.10", "memory" => 2048, "cpus" => 2 },
}
(1..NB_WORKERS).each do |n|
  MACHINES["worker#{n}"] = { "ip" => "192.168.56.#{10 + n}",
                             "memory" => 3072, "cpus" => 2 }
end
\end{lstlisting}

\begin{retenir}
Rendre le nombre de workers paramétrable coûte trois lignes et change tout
pour qui lit le projet : on peut le faire tourner sur un portable avec un
seul worker, ou monter à neuf pour une démonstration. Un fichier qui n'offre
qu'une seule taille finit toujours par être modifié à la main --- et donc
désynchronisé.
\end{retenir}

\section{Plusieurs providers, et pourquoi}

\begin{lstlisting}[language=Ruby]
node.vm.provider "vmware_desktop" do |v| ... end
node.vm.provider "parallels"      do |p| ... end
node.vm.provider "qemu"           do |q, override| ... end
node.vm.provider "virtualbox"     do |v| ... end
\end{lstlisting}

VirtualBox, le provider par défaut, ne fonctionne pas sur les Mac à puce
Apple. Déclarer plusieurs providers permet au même fichier de marcher sur une
machine Intel, un Mac ARM et un poste Windows.

\begin{piege}{QEMU ne gère pas les réseaux privés}
Chaque provider a ses limites. \cmd{vagrant-qemu} ne sait pas créer de réseau
privé : l'accès SSH est redirigé sur \cmd{localhost}. Le \cmd{Vagrantfile} le
détecte et ajuste la configuration. Sans cela, le déploiement réussirait et
rien ne serait joignable.
\end{piege}

\section{Les commandes du quotidien}

\begin{lstlisting}[language=bash]
vagrant up                  # cree et demarre tout le parc
vagrant up worker2          # une seule machine
vagrant status              # qui tourne
vagrant ssh controller      # entrer dans une machine
vagrant halt worker1        # eteindre (le test de panne du cahier des charges)
vagrant reload worker1      # redemarrer
vagrant destroy -f          # tout effacer, pour repartir de zero
\end{lstlisting}

\chapter{Configurer avec Ansible}
\label{chap:ansible}

\section{La partition : quatre plays, neuf rôles}

\begin{lstlisting}[language=yamlplus,caption={\fichier{3\_ansible/site.yml} --- la structure}]
- name: 1. Durcissement de tous les noeuds (UFW)
  hosts: insacloud
  become: true
  roles:
    - role: security_hardening

- name: 2. Controleur — compte de service et cle SSH de pilotage
  hosts: controllers
  roles:
    - role: controller_key

- name: 3. Workers — Docker, certificat, acces du controleur, images
  hosts: workers
  roles:
    - role: docker
    - role: tls_cert
    - role: worker
    - role: docker_images
    - role: worker_agent

- name: 4. Controleur — client Docker, certificat, nginx, application
  hosts: controllers
  roles:
    - role: docker
    - role: tls_cert
    - role: reverse_proxy
    - role: webapp
\end{lstlisting}

\begin{notion}{Pourquoi cet ordre}
L'ordre n'est pas arbitraire. Le contrôleur doit avoir sa clé SSH
\emph{avant} que les workers n'autorisent cette clé ; les workers doivent
avoir Docker \emph{avant} qu'on y construise les images ; nginx doit être
configuré \emph{avant} que l'application démarre derrière lui.

Un playbook est une suite de dépendances. Les écrire explicitement évite les
« ça marche une fois sur deux ».
\end{notion}

\begin{center}
\small
\begin{tabularx}{\textwidth}{@{}lX@{}}
\toprule
\textbf{Rôle} & \textbf{Ce qu'il fait} \\
\midrule
\cmd{security\_hardening} & pare-feu UFW, ports ouverts selon le groupe \\
\cmd{controller\_key} & compte de service et clé SSH de pilotage \\
\cmd{docker} & Docker Engine sur les workers, client seul sur le contrôleur \\
\cmd{tls\_cert} & certificat auto-signé propre à chaque nœud \\
\cmd{worker} & accès SSH restreint du contrôleur, appartenance au groupe docker \\
\cmd{docker\_images} & construction des images des machines louées \\
\cmd{worker\_agent} & l'agent d'inscription et de battement de cœur \\
\cmd{reverse\_proxy} & nginx HTTPS devant Gunicorn \\
\cmd{webapp} & Flask, Gunicorn, Faucheur, watchdog, secrets \\
\bottomrule
\end{tabularx}
\end{center}

\section{Un piège réel : la déduplication des rôles}

Ansible n'exécute un rôle qu'une fois par play. Sauf si on l'appelle avec des
\emph{paramètres différents} : il le considère alors comme deux rôles
distincts et l'exécute deux fois.

Le rôle \cmd{docker} était appelé directement par le play du contrôleur, et
tiré en dépendance par \cmd{docker\_images}. Avec un paramètre différent à
chaque appel, il s'exécutait deux fois --- installant puis désinstallant le
démon.

\begin{lstlisting}[language=yamlplus,caption={La correction : une variable de play, pas un paramètre de rôle}]
- name: 4. Controleur
  hosts: controllers
  vars:
    single_host: "{{ groups['workers'] | default([]) | length == 0 }}"
    # Defini ici (et non en parametre de role) : un role appele avec des
    # parametres differents n'est pas dedoublonne par Ansible, et `docker` est
    # aussi tire comme dependance par `docker_images`. Une variable de play
    # s'applique aux deux appels, qui sont donc fusionnes en une execution.
    docker_cli_only: "{{ not single_host }}"
  roles:
    - role: docker
    - role: docker_images
      when: single_host
\end{lstlisting}

\begin{retenir}
Quand un comportement d'outil vous surprend, vérifiez-le par l'expérience
plutôt que par la documentation. Nous avons prouvé celui-ci en instrumentant
le rôle pour qu'il journalise chaque exécution.
\end{retenir}

\section{Les secrets : Ansible Vault}

\begin{lstlisting}[language=bash]
# Creer ou editer un fichier chiffre
ansible-vault create group_vars/insacloud/vault.yml
ansible-vault edit   group_vars/insacloud/vault.yml

# Jouer le playbook en fournissant la phrase de passe
ansible-playbook site.yml --ask-vault-pass
\end{lstlisting}

Le fichier est chiffré en AES-256 et peut être versionné sans danger. Il
contient la clé de session Flask, la clé du coffre Fernet, le mot de passe
PostgreSQL et le jeton des agents.

\begin{piege}{Un secret déposé dans une unité systemd}
Mettre \cmd{Environment=INSACLOUD\_VAULT\_KEY=...} dans un fichier d'unité le
rend lisible par n'importe qui : les unités sont en lecture pour tous, et
\cmd{systemctl show} les affiche.

La bonne façon : un fichier séparé en \cmd{0600}, appartenant au compte de
service, référencé par \cmd{EnvironmentFile=}. Le secret n'est alors ni dans
l'unité, ni dans les journaux.
\end{piege}

\section{Prouver l'idempotence}

\begin{lstlisting}[language=bash]
ansible-playbook site.yml --ask-vault-pass        # premier passage
ansible-playbook site.yml --ask-vault-pass        # second passage
\end{lstlisting}

Le second passage doit afficher \cmd{changed=0} sur chaque nœud. C'est mesuré
sur les quatre machines après chaque modification du projet.

\begin{center}
\small
\begin{verbatim}
PLAY RECAP ****************************************************************
controller : ok=71  changed=0  unreachable=0  failed=0  skipped=9
worker1    : ok=43  changed=0  unreachable=0  failed=0  skipped=4
worker2    : ok=43  changed=0  unreachable=0  failed=0  skipped=4
worker3    : ok=43  changed=0  unreachable=0  failed=0  skipped=4
\end{verbatim}
\end{center}

\begin{piege}{\texttt{command} et \texttt{shell} ne sont jamais idempotents}
Les modules \cmd{ansible.builtin.command} et \cmd{shell} exécutent ce qu'on
leur donne, sans rien vérifier : ils sont donc \emph{toujours} signalés
\cmd{changed}. Si vous devez les employer, ajoutez \cmd{creates:},
\cmd{removes:} ou \cmd{changed\_when:} pour leur apprendre à se taire quand il
n'y a rien à faire.
\end{piege}

\section{L'inventaire dynamique}

L'inventaire statique liste les machines à la main. Un inventaire
\emph{dynamique} les découvre.

\begin{lstlisting}[language=bash]
./inventaire_vagrant.py --list      # interroge Vagrant, produit du JSON
ansible-playbook -i inventaire_vagrant.py site.yml
\end{lstlisting}

Le script interroge \cmd{vagrant status} et construit les groupes. Changer le
nombre de VM ne demande plus de toucher à l'inventaire, et le déploiement
devient indépendant du provider de virtualisation.

\chapter{Servir l'application}

\section{nginx devant Gunicorn}

\begin{lstlisting}[language=bash,caption={Le principe de la configuration nginx}]
server {
    listen 443 ssl http2;
    ssl_protocols TLSv1.3;

    add_header Strict-Transport-Security "max-age=63072000" always;
    add_header X-Content-Type-Options nosniff always;

    location /static/ {
        alias /opt/insacloud/static/;    # servi par nginx, pas par Python
        expires 7d;
    }

    location / {
        proxy_pass http://127.0.0.1:5000;
        proxy_set_header Host $host;
        proxy_set_header X-Forwarded-For   $proxy_add_x_forwarded_for;
        proxy_set_header X-Forwarded-Proto $scheme;
    }
}
\end{lstlisting}

\begin{notion}{\texttt{X-Forwarded-For} et \texttt{ProxyFix}}
Derrière un proxy, toutes les requêtes semblent venir de \cmd{127.0.0.1}.
L'adresse réelle du client est placée par nginx dans l'en-tête
\cmd{X-Forwarded-For}.

Côté Flask, le \emph{middleware} \cmd{ProxyFix} lit ces en-têtes et rétablit
la bonne adresse. Sans lui, notre verrouillage anti-force-brute compterait
toutes les tentatives comme venant de la même IP --- et bloquerait tout le
monde au cinquième échec de n'importe qui.
\end{notion}

\section{systemd : faire tourner les services}

\begin{lstlisting}[caption={Le principe d'une unité systemd}]
[Unit]
Description=InsaCloud - application web
After=network-online.target

[Service]
User=insacloud
WorkingDirectory=/opt/insacloud
EnvironmentFile=/etc/insacloud/insacloud.env     # secrets, en 0600
ExecStart=/opt/insacloud/.venv/bin/gunicorn --preload -w 4 -b 127.0.0.1:5000 app:app
Restart=always
RestartSec=3

[Install]
WantedBy=multi-user.target
\end{lstlisting}

\begin{center}
\small
\begin{tabularx}{\textwidth}{@{}lX@{}}
\toprule
\textbf{Directive} & \textbf{Effet} \\
\midrule
\cmd{User=} & le service ne tourne pas en root \\
\cmd{EnvironmentFile=} & les secrets viennent d'un fichier protégé, pas de l'unité \\
\cmd{Restart=always} & systemd relance le service s'il meurt \\
\cmd{After=network-online} & ne pas démarrer avant que le réseau soit prêt \\
\cmd{--preload} & Gunicorn importe l'application une fois, puis duplique \\
\bottomrule
\end{tabularx}
\end{center}

\section{Le watchdog}

Un \emph{timer} systemd vérifie chaque minute que nginx répond, que le site
renvoie bien du HTTPS, et que le Faucheur tourne. Il redémarre ce qui est
tombé et journalise les workers injoignables.

\begin{notion}{Timer plutôt que cron}
Un timer systemd s'intègre aux journaux (\cmd{journalctl -u ...}), connaît
l'état du service qu'il surveille, et ne démarre pas avant que ses
dépendances soient prêtes. Cron ne sait rien de tout cela.
\end{notion}

\chapter{La pile \texttt{docker compose}}

\section{Pourquoi une seconde façon de déployer}

Monter quatre VM demande du temps, de la mémoire et un hyperviseur. Pour
développer, démontrer, ou faire tourner les tests d'intégration en
intégration continue, on veut une pile en une commande.

\begin{lstlisting}[language=bash]
cp .env.example .env         # puis renseigner les secrets
docker compose up -d
\end{lstlisting}

\section{Un bloc partagé par deux services}

\begin{lstlisting}[language=yamlplus,caption={\fichier{docker-compose.yml} --- l'ancre YAML}]
x-application: &application
  build:
    context: ./2_webapp
  image: insacloud_web:latest
  restart: unless-stopped
  environment:
    INSACLOUD_ENV: production
    INSACLOUD_CLUSTER: ${INSACLOUD_CLUSTER:-compose}
    DATABASE_URL: "postgresql+psycopg://...:${POSTGRES_PASSWORD:?manquant}@db:5432/..."
    INSACLOUD_SECRET_KEY: ${INSACLOUD_SECRET_KEY:?definissez INSACLOUD_SECRET_KEY}
  depends_on:
    db:
      condition: service_healthy

services:
  web:
    <<: *application              # reprend tout le bloc ci-dessus
    ports: ["127.0.0.1:8088:5000"]
  faucheur:
    <<: *application
    command: ["python", "faucheur.py"]
\end{lstlisting}

\begin{notion}{Les ancres YAML}
\cmd{\&application} donne un nom à un bloc ; \cmd{<<: *application} le
recopie. L'application et le Faucheur partagent ainsi exactement la même
configuration --- il devient impossible qu'ils divergent sur une variable.
\end{notion}

\begin{notion}{\texttt{\$\{VAR:?message\}}}
Cette syntaxe refuse de démarrer si la variable n'est pas définie, en
affichant le message. La pile ne peut donc pas tourner avec un mot de passe
vide ou une clé par défaut : l'erreur arrive au démarrage, pas trois semaines
plus tard.
\end{notion}

\section{Le volume partagé, et pourquoi il ne s'invente pas}

\begin{lstlisting}[language=yamlplus]
  web:
    volumes:
      - /var/run/docker.sock:/var/run/docker.sock
      - archives_travail:/var/lib/insacloud/archives
  faucheur:
    volumes:
      - /var/run/docker.sock:/var/run/docker.sock
      - archives_travail:/var/lib/insacloud/archives

volumes:
  donnees_postgres:
  archives_travail:
\end{lstlisting}

\begin{piege}{Deux services, un seul dossier}
Le Faucheur \emph{écrit} l'archive ; l'application la \emph{sert} au
téléchargement. Si chacun avait son propre volume, l'utilisateur verrait dans
son tableau de bord une archive que le serveur ne trouverait pas.

Le même volume nommé doit donc être monté dans les deux services. C'est le
genre d'erreur qui ne se voit pas à la lecture du fichier, et qui saute aux
yeux au premier essai réel.
\end{piege}

\begin{piege}{La socket Docket donne les pleins pouvoirs}
Monter \cmd{/var/run/docker.sock} dans un conteneur lui donne un pouvoir
équivalent à root sur l'hôte : il peut créer un conteneur privilégié qui monte
le disque entier. C'est acceptable sur un poste de développement, jamais en
production --- où le contrôleur pilote des workers distants par SSH et n'a
plus besoin de cette ligne.
\end{piege}

\section{Les profils}

\begin{lstlisting}[language=bash]
docker compose up -d                        # la plateforme seule
docker compose --profile build build        # + construire les six images
docker compose --profile demo up -d         # + une machine de demonstration
\end{lstlisting}

Un service marqué d'un profil ne démarre que si on le demande. Cela permet de
décrire dans un seul fichier la pile minimale et tous ses compléments.
